% Layout fixture: displays that TeX sets as one printed line or one printed equation but that
% pieces of the page can make look like several: two formulas joined by "and" or "for all", and
% a long aligned row whose right half starts clear of the row above it. Used by
% tests/test_layout_display_groups.py.
% Written for ReadableSTEM's tests; no third-party text.
\documentclass[11pt]{article}
% Computer Modern text fonts (CMR, CMBX), as pdfLaTeX writes them; XeTeX defaults to Latin Modern.
\usepackage[OT1]{fontenc}
\renewcommand{\rmdefault}{cmr}
\usepackage{amsmath,amssymb}
\usepackage[margin=1in]{geometry}
\pagestyle{empty}

\begin{document}

\section*{Two formulas on one line}

The two partition coefficients of a species are the ratios of its concentrations in the
phases,
\begin{equation}
  P_A = \frac{c_A^{(1)}}{c_A^{(2)}} \hspace{5em} \text{and} \hspace{5em}
  P_B = \frac{c_B^{(1)}}{c_B^{(2)}} \;,
\end{equation}
which can be measured. In a trial, suppose that
\[
  P(A = 1 \mid L = 1) = 0.75 \quad \text{and} \quad P(A = 1 \mid L = 0) = 0.50.
\]
The weights must satisfy
\[
  0 < \pi(L_i) < 1 \quad \text{for all } i.
\]

\section*{A long aligned row}

Shifting the contour gives
\begin{align*}
  \mathbb{E}[L_n] &= -\frac{1}{2\pi \mathrm{i}} \int_{\sigma-\mathrm{i}\infty}^{\sigma+\mathrm{i}\infty}
    \frac{\widetilde{f}_n(s)}{s}\,\mathrm{d}s \\
  &= -\frac{1}{2\pi \mathrm{i}} \int_{\sigma-\mathrm{i}\infty}^{\sigma+\mathrm{i}\infty}
    \frac{1}{s^2(\psi(1-s)-\psi(1))}\,\mathrm{d}s
    + \frac{1}{2\pi \mathrm{i}} \int_{\sigma-\mathrm{i}\infty}^{\sigma+\mathrm{i}\infty}
    \frac{\Gamma(s)\Gamma(n)/\Gamma(n+s)}{s(\psi(1-s)-\psi(1))}\,\mathrm{d}s.
\end{align*}

\section*{Lists that go on}

The gaps are computed for $j = 1,2,3,...$ obtained by diagonalisation, and the roots
are ordered as $0 > s_1 > s_2 > ...$ denoting the negative roots.

\section*{Words that do not join}

We have $x = 1$ and $y = 2$ in this case, and nothing else. When
\[
  a = b \qquad \text{where we assume} \qquad c = d,
\]
the claim holds. Two numbered equations stay apart,
\begin{align}
  u &= \frac{1}{2} \\
  v &= \frac{1}{3}
\end{align}
and so do two unrelated lines.

\newpage
\noindent $p_1 = 3$ and $q_1 = 4$\\
is the tail of a sentence from the page before: it starts at the left margin and is no display.
The rest of this page is running text, so that the column's margins can be measured against
it, and it runs on for a few lines across the whole width of the column in the usual way.

\end{document}
